秋天的时候，距离我们那次地铁偶遇已经过去了半年。我已经记不清确切的日子，只记得第一次约见是在北海公园。刚下过一场小雨，石板路上泛着薄薄的水渍。她穿了一双白色的帆布鞋，走了一段，鞋尖很快就被泥水洇湿了，但她低头瞥了一眼，似乎并不在意。

这些细节我在当时并不太关注，直到事后才逐渐回想起来。想起半年来的恋爱，不，与其说是恋爱，倒不如说是我的一厢情愿，我们并没在事实上达成一种心照不宣的默契，大部分的情感流动，更像是我单方面的揣度。那时的我对理佳有一种病态般的迷恋，我时时给她发一些无关痛痒的消息，以显示我对她的关注，或者说，如果我无法与她构成互动，我害怕她会如烟一般从世上消失。

走到湖边的时候，她突然停下来蹲下去，角落里蹲着一只橘猫，正眯着眼睛晒太阳。她伸出手，猫犹豫了一下，把头蹭过来，我在她身后跟着蹲下去，她的头发垂下来正好散在我的眼前，很近。我闻到了她头发上的味道，混在雨水后空气里的湿润，说不出是什么，但很好闻。她侧过脸和猫嬉戏，没有注意到我。我突然有一种急剧的冲动，于是将脸更加凑近了一些。她回头看了我一眼，并没有说话。

她注意到我了吗？她会怎么看待我呢？她会将我看成轻浮而不成熟的人吗？不，我就应该勇敢一些，这有什么关系，难道情感的诞生不是因为行为的越界吗？如果一直保持谦恭有礼的模样，那么关系该怎么更进一步呢？

我们走着走着走上来了白塔，下来的时候路很陡，她走在前面，非常小心。我走在后面，突然跑了两步，从最后几级台阶上一蹦而下。落地的时候声音很大，在塔底下荡开来，她十分惊讶地看着我，忽而笑出了声。

“你怎么也这样呀”

她笑得太厉害了，甚至眼角都渗出一点湿意，我也笑了，在那一刻我仿佛又回到了小时候，并非是在读中学的时候，而是在更小的时候那种无忧无虑的感觉。即使现在想起，我的嘴角依然会不由自主地浮起笑容。

“也”吗？这个也包含谁呢？我突然又被一阵现实撕扯回来。我原本以为我是在做一件独一无二的事，然而在她充满笑意的眼睛中，我突然感受到微微的失落。

那天晚上，我们在胡同里的一家居酒屋吃饭。几杯清酒下肚后，也许是灯光太暗，她突然看着我，然后指向居酒屋角落里的一把旧木吉他。

“说起来，我从前也跟风学过一段时间的吉他，当时的吉他老师问我，你知道吉他为什么会发出这么好听的共鸣吗？”

“因为琴箱？” 我顺着她的目光看过去，回答道。

“对，因为琴箱是空的。”她收回目光，看着自己面前的空酒杯，“如果那是一块实心的木头，你拨动琴弦，它只会发出沉闷的死声音。它之所以好听与迷人，恰恰因为它什么都没有，它只是在放大你拨弦时的情绪。”

理佳沉吟了一会，用空酒杯敲了敲桌子，

“我记得吉他老师当时一再告诫我，扫弦的时候千万不要太用力。 因为琴箱是空的，面板很脆弱，承受不了太重的敲击。如果弹得太狠，它是会裂开的。”

“有点像在谈禅了。”我说道。

“倒也没有那么深奥，只是看到了这有一把吉他，突然想到。” 她又恢复了以往的笑容。

“不过，你学吉他这件事，是从什么时候开始的，怎么从来没有听过。”

“在本科的时候参加社团，出于好奇学了的吉他，不过再参加完一届社团的表演后就再也没有碰过了。”

“后面为什么没有坚持弹下去呢。”

“后面退出了社团，身边也没有人弹了，就逐渐搁置了。”

我注意到，她去做一件事或者放弃一件事总是和人相关，她总是在说一些客观的原因与现实，却很少去表达自己对某件事的深刻感受。

沉默有顷。

我喝完了杯中最后一杯酒，看了眼手表，时间已经九点了，“走吧，”，我说，“时间不早了。”

“嗯。”

我们起身，临出门前，她突然问我，你喜欢猫吗？

“喜欢倒是喜欢，可要是我养的话，却总觉得麻烦，毕竟我连自己都顾不好呢。”

她没有说话，似乎只是将这句话当成了一句普通的感慨，我却总觉得心里坠着什么，这句话似乎比我想象的泄露了更多，但是具体泄露了多少，会造成什么影响，我自己也说不清楚。

十一月的一个晚上，我给她打电话。响了很久，没有人接。我挂掉，又打了一次。响了几声，断了。我再打。关机。

之后我再也没有打通过那个号码。

十二月初我去了她的学校。传达室的人说，好像有这么个人，去年就离职了。我问去了哪里。不知道。我站在校门口，风从衣领往里灌。门口有几个学生蹲在路边吃烤串，烟气被风吹散了。我忽然想起来，我不知道她教什么科目。

和预感的一样，理佳果真如烟一般消失了。

第二天傍晚，我希望能够去理佳常去的地铁站碰碰运气，或许可以碰到理佳。我站在四号线与二号线的交叉的拐角长廊，在汹涌的、面目模糊的人潮中站了半个小时。每一个身着短上衣或者走路姿势有些像她的人路过，我的心跳都会沉重地错漏一拍，直到值班的站务员用狐疑的目光打量我，我才意识到自己的荒诞，狼狈地顺着电梯出了站。

我站在出口的冷风中，手机突然震动，我收到了一个陌生号码发送来的短信，里面写道，

屏幕上只有一行字，“亮，请原谅我的离开，也许是我太冷漠了，但请不要再找我。”

我几乎是本能地把电话拨了过去，耳筒里传来的不是忙音，而是一段毫无感情的机械女声：“您好，您拨打的号码是空号……”

我又拨了一次，依然是那段机械的提示音。那一刻，地铁隆隆驶过地底的震动顺着我的鞋底传上来。
